\documentclass[11pt,a4paper]{report}
\usepackage[margin=2.5cm]{geometry}
\usepackage{graphicx}
\usepackage{hyperref}
\usepackage{longtable}
\usepackage{array}
\usepackage{titlesec}
\usepackage{enumitem}

\titleformat{\chapter}{\normalfont\huge\bfseries}{\thechapter.}{1em}{}

\begin{document}

\begin{titlepage}
    \centering
    \vspace*{1.5cm}
    {\LARGE Politecnico di Milano\par}
    \vspace{0.5cm}
    {\Large Stakeholder Expense Report\par}
    \vspace{1.5cm}
    {\huge BUDD-e LiDAR Pedestrian Annotation Pipeline\par}
    \vspace{1.0cm}
    {\large Project Period: \underline{\hspace{5cm}}\par}
    \vspace{0.5cm}
    {\large Prepared by: \underline{\hspace{7cm}}\par}
    \vspace{0.5cm}
    {\large Contact: \underline{\hspace{8cm}}\par}
    \vfill
    {\large Date: \underline{\hspace{5cm}}\par}
\end{titlepage}

\tableofcontents
\clearpage

\chapter{Executive Summary}
This report documents the work completed to establish a LiDAR-only pipeline for pedestrian 3D bounding box annotation from ROS1 bag data recorded by the BUDD-e mobile robot in a hospital environment. The work includes data extraction, model inference with OpenPCDet, export of labels, and integration with Segments.ai for review and annotation.

\chapter{Objectives}
\begin{itemize}[leftmargin=2em]
    \item Enable ROS bag ingestion and conversion to OpenPCDet-compatible point cloud frames.
    \item Run pretrained LiDAR-only 3D detection to generate pedestrian bounding boxes.
    \item Export annotations in reusable formats and upload to a web-based labeling UI.
    \item Prepare the workflow for future tracking/association across frames.
\end{itemize}

\chapter{Data and Inputs}
\begin{itemize}[leftmargin=2em]
    \item ROS1 bag files containing \texttt{/rslidar\_points} (sensor\_msgs/PointCloud2).
    \item Example bag used for development: \texttt{File\_2023-06-27-16-15-39.bag} (2:07 min).
    \item Generated point clouds: \texttt{.bin} and \texttt{.pcd} sequence files.
\end{itemize}

\chapter{Pipeline Overview}
\section*{Ingestion}
\begin{itemize}[leftmargin=2em]
    \item ROS bag inspection to verify topics and message counts.
    \item Extraction of \texttt{/rslidar\_points} into per-frame point clouds.
    \item Conversion to OpenPCDet format: \texttt{(x, y, z, intensity)} float32.
\end{itemize}

\section*{Detection}
\begin{itemize}[leftmargin=2em]
    \item OpenPCDet installed and configured with CUDA-enabled PyTorch.
    \item Baseline model: PointPillars (KITTI pretrained).
    \item Outputs: 3D bounding boxes and confidence scores per frame.
\end{itemize}

\section*{Label Export and Review}
\begin{itemize}[leftmargin=2em]
    \item Labels exported to JSONL and MATLAB formats.
    \item PCD sequences prepared for external annotation tools.
    \item Segments.ai dataset created for 3D web-based review and correction.
\end{itemize}

\chapter{Work Performed}
\begin{itemize}[leftmargin=2em]
    \item Installed and configured ROS Noetic in WSL.
    \item Installed OpenPCDet dependencies, CUDA bindings, and compiled ops.
    \item Implemented scripts:
    \begin{itemize}[leftmargin=2em]
        \item ROS bag to OpenPCDet frames extraction.
        \item Model inference over frame sequences.
        \item Export to MATLAB and PCD sequences.
        \item Segments.ai upload of sequences and prelabels.
    \end{itemize}
    \item Uploaded prelabels to Segments.ai (pointcloud cuboid sequence).
\end{itemize}

\chapter{Deliverables}
\begin{longtable}{p{0.35\textwidth} p{0.6\textwidth}}
\textbf{Deliverable} & \textbf{Description} \\
\hline
Point cloud frames & OpenPCDet-compatible \texttt{.bin} files per frame \\
PCD sequence & \texttt{.pcd} sequence for external visualization tools \\
Predictions JSONL & Raw model outputs (boxes, scores, labels) \\
Pedestrian export & Filtered pedestrian-only labels (JSONL/CSV/MAT) \\
Segments.ai dataset & Web UI with prelabels for review and correction \\
\end{longtable}

\chapter{Issues and Mitigations}
\begin{itemize}[leftmargin=2em]
    \item \textbf{Model performance}: Baseline KITTI pretrained PointPillars is weak for indoor pedestrians.
    \item \textbf{Viewer performance}: Large sequences can overload web labeling UI.
    \item \textbf{Mitigation}: Split sequences into smaller chunks and use pedestrian-only filtering for review.
\end{itemize}

\chapter{Next Steps}
\begin{itemize}[leftmargin=2em]
    \item Evaluate stronger LiDAR-only models (CenterPoint, PV-RCNN).
    \item Apply simple tracking to associate boxes across frames.
    \item Refine upload pipeline with smaller clips and higher-confidence prelabels.
    \item Automate end-to-end script with arguments:
    \begin{itemize}[leftmargin=2em]
        \item \texttt{--bag}, \texttt{--topic}, \texttt{--model-config}, \texttt{--checkpoint}
        \item \texttt{--labelset}, \texttt{--upload}, \texttt{--chunk-size}
    \end{itemize}
\end{itemize}

\chapter{Expense Summary}
Please complete this section with the project-specific costs and time spent.
\begin{itemize}[leftmargin=2em]
    \item Personnel hours: \underline{\hspace{5cm}}
    \item Compute resources: \underline{\hspace{6cm}}
    \item Software/licenses: \underline{\hspace{6cm}}
    \item Other expenses: \underline{\hspace{6cm}}
\end{itemize}

\end{document}
